% Part: incompleteness
% Chapter: theories-computability
% Section: q-is-ce

\documentclass[../../../include/open-logic-section]{subfiles}

\begin{document}

\olfileid{inc}{tcp}{qce}

\olsection{$\Th{Q}$ એ \printtoken{S}{c.e.}-પૂર્ણ છે}


\begin{thm}
$\Th{Q}$ !!{c.e.} છે,
પણ નિર્ણેય નથી. હકીકતમાં તે
પૂર્ણ !!{c.e.} ગણ છે.
\end{thm}

\begin{proof}
$\Th{Q}$ !!{c.e.} છે
તે જોવું મુશ્કેલ નથી, કારણ કે
તે એવા વાક્યોના (સંકેતોના)
ગણ~$y$ છે કે $\Th{Q}$માં
$y$ની કોઈ સાબિતી~$x$ છે:
\[
Q = \Setabs{y}{\lexists[x][\Prf[\Th{Q}](x,y)]}.
\]
પરંતુ આપણે જાણીએ છીએ કે
$\Prf[\Th{Q}](x,y)$ સંગણનીય
છે (હકીકતમાં આદિમ પુનરાવર્તી),
અને ઉપરના આકારમાં લખી શકાય
એવો દરેક ગણ સંગણકીય રીતે
પરિગણનીય છે.

તે પૂર્ણ સંગણકીય રીતે
પરિગણનીય ગણ છે એમ કહેવું
$K \leq_m Q$ કહેવાની સમકક્ષ
છે, જ્યાં $K = \Setabs{x}{!A_x(x) \downarrow}$.
તેથી $K$નું $\Th{Q}$માં
ન્યૂનીકરણ થાય છે તે બતાવીએ.
ક્લીનીનું વિધેયધર્મ
$T(e,x,s)$ આદિમ પુનરાવર્તી
હોવાથી તે $\Th{Q}$માં,
કહો કે~$!A_T$ વડે, નિરૂપ્ય
છે. તેથી દરેક~$x$ માટે
\begin{align*}
x \in K & \lif \lexists[s][T(x,x,s)] \\
& \lif \lexists[s][(\Th{Q} \Proves !A_T(\num x,\num x, \num s))]
  \\
& \lif \Th{Q} \Proves \lexists[s][!A_T(\num x, \num x, s)].
\end{align*}
વિપરીત દિશામાં, જો
$\Th{Q} \Proves \lexists[s][!A_T(\num x, \num x, s)]$,
તો હકીકતમાં કોઈ પ્રાકૃતિક
સંખ્યા~$n$ માટે સૂત્ર
$!A_T(\num x, \num x, \num n)$
સાચું હોવું જ પડે. હવે
$T(x,x,n)$ ખોટું હોય, તો
$!A_T$ એ~$T$ને નિરૂપે છે
એટલે $\Th{Q}$માં
$\lnot !A_T(\num x, \num x, \num n)$
સાબિત થાય. પણ ત્યારે
$\Th{Q}$માં ખોટું સૂત્ર
સાબિત થતું હોત; આ
વિરોધાભાસ છે. તેથી
$T(x,x,n)$ સાચું જ હોવું
જોઈએ, અને પરિણામે
$!A_x(x) \downarrow$.

ટૂંકમાં, દરેક~$x$ માટે
$x$ એ~$K$માં છે ત્યારે અને
માત્ર ત્યારે જ $\Th{Q}$માં
$\lexists[s][!A_T(\num x,\num x,s)]$
સાબિત થાય છે. તેથી~$x$ને
વાક્ય $\lexists[s][!A_T(\num x,
  \num x, s)]$નો સંકેત આપતું
વિધેય~$f$ એ~$K$નું
$\Th{Q}$માં ન્યૂનીકરણ છે.
\sourcecorrection{OLINC-039}{સ્થિર અંગ્રેજી
મૂળ આ અંતિમ બે વાક્યોમાં
અંકગણિતની ભાષામાં વપરાતું
નિરૂપક સૂત્ર~$!A_T$ બદલે
બાહ્ય ક્લીની વિધેયધર્મ~$T$
સૂત્રની અંદર લખે છે.
ન્યૂનીકરણનું વાક્ય અગાઉની
સાબિતી પ્રમાણે $!A_T$
વડે જ રચ્યું છે.}
\end{proof}

\end{document}
